\chapter{压杆稳定性}

\section{压杆稳定的基本概念}

\begin{definition}[title=失稳与临界压力]
对于细长压杆，当轴向压力 \(F\) 小于某一极限值时，杆件保持直线平衡状态；若受到微小侧向干扰，撤去干扰后杆件能恢复直线平衡，则称该平衡是\textbf{稳定的}。当压力 \(F\) 超过该极限值时，微小的干扰会使杆件发生显著弯曲变形甚至破坏，这种平衡状态称为\textbf{不稳定的}。压力由稳定平衡过渡到不稳定平衡的临界值称为\textbf{临界压力}，记作 \(F_{\text{cr}}\)。
\end{definition}

\begin{theorem}[title=临界压力的实验现象]
当作用在细长压杆上的轴向压力达到临界压力 \(F_{\text{cr}}\) 时，杆件可以保持在微弯状态下的平衡；压力稍大于 \(F_{\text{cr}}\) 时，弯曲变形迅速增大，最终导致压杆失效。
\end{theorem}

\section{细长压杆的临界压力——欧拉公式}

考虑两端铰支的理想细长压杆，长度为 \(l\)，弯曲刚度为 \(EI\)。设压力 \(F\) 达到临界值 \(F_{\text{cr}}\) 时杆件处于微弯平衡状态，挠曲线方程为 \(w(x)\)。由弯矩平衡得：
\[
M(x) = -F_{\text{cr}} w(x)
\]
而梁的近似挠曲线微分方程为：
\[
EI w''(x) = M(x) = -F_{\text{cr}} w(x)
\]
令 \(k^2 = \dfrac{F_{\text{cr}}}{EI}\)，则：
\[
w'' + k^2 w = 0
\]
通解为：
\[
w(x) = A \sin kx + B \cos kx
\]
利用边界条件 \(w(0)=0\) 得 \(B=0\)；再由 \(w(l)=0\) 得 \(A \sin kl = 0\)。若 \(A \neq 0\)（非平凡解），则需 \(\sin kl = 0\)，故：
\[
kl = n\pi \quad (n=1,2,3,\dots)
\]
取最小非零解 \(n=1\)，得到临界压力：
\[
F_{\text{cr}} = \frac{\pi^2 EI}{l^2}
\]
此为两端铰支压杆的\textbf{欧拉公式}。

\begin{theorem}[title=欧拉公式（两端铰支）]
两端铰支的细长压杆，其临界压力为：
\[
\boxed{F_{\text{cr}} = \frac{\pi^2 EI}{l^2}}
\]
式中 \(E\) 为材料的弹性模量，\(I\) 为横截面的最小形心主惯性矩，\(l\) 为压杆长度。
\end{theorem}

\section{不同约束条件下的欧拉公式}

对于不同支座条件，只需将欧拉公式中的长度 \(l\) 换算为\textbf{相当长度} \(\mu l\)，其中 \(\mu\) 称为\textbf{长度因数}。临界压力的一般形式为：
\[
F_{\text{cr}} = \frac{\pi^2 EI}{(\mu l)^2}
\]
常见约束的长度因数：
\begin{center}
\begin{tabular}{|c|c|c|}
\hline
约束形式 & 长度因数 \(\mu\) & 失稳时挠曲线形状特征 \\ \hline
两端铰支 & \(1.0\) & 正弦半波 \\ \hline
一端固定、一端自由 & \(2.0\) & 正弦四分之一波 \\ \hline
一端固定、一端铰支 & \(0.7\) & 正弦0.7波 \\ \hline
两端固定 & \(0.5\) & 正弦半波（两端拐点） \\ \hline
\end{tabular}
\end{center}

\begin{proposition}[title=柔度与临界应力]
压杆的柔度（长细比）定义为：
\[
\lambda = \frac{\mu l}{i}, \quad i = \sqrt{\frac{I}{A}}
\]
其中 \(i\) 为截面的惯性半径。临界应力为：
\[
\sigma_{\text{cr}} = \frac{F_{\text{cr}}}{A} = \frac{\pi^2 E}{\lambda^2}
\]
\end{proposition}

\section{欧拉公式的适用范围}

欧拉公式推导中使用了线弹性小变形假设，因此临界应力不能超过材料的比例极限 \(\sigma_p\)：
\[
\frac{\pi^2 E}{\lambda^2} \le \sigma_p \quad \Rightarrow \quad \lambda \ge \pi \sqrt{\frac{E}{\sigma_p}} \triangleq \lambda_p
\]
\(\lambda_p\) 称为\textbf{比例极限柔度}。当 \(\lambda \ge \lambda_p\) 时，欧拉公式适用，此类压杆称为\textbf{大柔度杆}（细长杆）。

常见材料的 \(\lambda_p\) 值：
\begin{itemize}
\item Q235钢：\(E \approx 206 \,\text{GPa},\ \sigma_p \approx 200 \,\text{MPa}\)，\(\lambda_p \approx 100\)
\item 铸铁：\(\lambda_p \approx 80\)
\item 木材：\(\lambda_p \approx 110\)
\end{itemize}

\section{中、小柔度杆的经验公式}

当 \(\lambda < \lambda_p\) 时，欧拉公式失效，需采用经验公式。

\begin{definition}[title=直线公式]
对于中柔度杆（\(\lambda \le \lambda_c\)，\(\lambda_c\) 为材料常数），临界应力与柔度成线性关系：
\[
\sigma_{\text{cr}} = a - b\lambda
\]
式中 \(a, b\) 为与材料有关的常数，可通过实验测定。例如 Q235钢：\(a = 304 \,\text{MPa}\)，\(b = 1.12 \,\text{MPa}\)。
\end{definition}

\begin{definition}[title=抛物线公式]
某些规范推荐采用抛物线公式：
\[
\sigma_{\text{cr}} = \sigma_s \left[ 1 - \alpha \left( \frac{\lambda}{\lambda_c} \right)^2 \right]
\]
其中 \(\sigma_s\) 为屈服极限，\(\alpha\) 为系数，\(\lambda_c\) 为抛物线公式适用的上限柔度。
\end{definition}

对于小柔度杆（\(\lambda\) 很小），压杆破坏主要是强度问题，取 \(\sigma_{\text{cr}} = \sigma_s\)（塑性材料）或 \(\sigma_{\text{cr}} = \sigma_b\)（脆性材料）。

\section{压杆的稳定校核}

为保证压杆安全工作，必须满足稳定条件：
\[
\frac{F_{\text{cr}}}{F} \ge n_{\text{st}} \quad \text{或} \quad \sigma_{\text{cr}} \ge \frac{F}{A} n_{\text{st}}
\]
其中 \(n_{\text{st}}\) 为规定的稳定安全系数，一般大于强度安全系数。工程上常采用\textbf{折减系数法}：
\[
\sigma = \frac{F}{A} \le \varphi [\sigma]
\]
式中 \(\varphi(\lambda)\) 为稳定系数（小于1），\([\sigma]\) 为材料的许用压应力。

\begin{example}[title=两端铰支钢压杆的临界压力]
已知 Q235 钢制压杆，截面为 \(50 \times 50 \,\text{mm}\) 的正方形，长度 \(l = 2 \,\text{m}\)，两端铰支，弹性模量 \(E = 206 \,\text{GPa}\)，比例极限 \(\sigma_p = 200 \,\text{MPa}\)。求临界压力并判断是否适用欧拉公式。
\end{example}
\begin{solution}
截面惯性矩：
\[
I = \frac{50^4}{12} = 520833 \,\text{mm}^4 = 5.208 \times 10^{-7} \,\text{m}^4
\]
两端铰支 \(\mu=1\)，欧拉公式给出：
\[
F_{\text{cr}} = \frac{\pi^2 \times 206 \times 10^9 \times 5.208 \times 10^{-7}}{2^2}
= \frac{\pi^2 \times 206 \times 5.208 \times 10^2}{4} \approx 2.646 \times 10^5 \,\text{N} = 264.6 \,\text{kN}
\]
临界应力：
\[
\sigma_{\text{cr}} = \frac{F_{\text{cr}}}{A} = \frac{264.6 \times 10^3}{0.05^2} = 105.8 \,\text{MPa}
\]
柔度：
\[
i = \sqrt{\frac{I}{A}} = \sqrt{\frac{5.208\times10^{-7}}{0.0025}} = 0.01443 \,\text{m},\quad
\lambda = \frac{\mu l}{i} = \frac{1\times 2}{0.01443} \approx 138.6
\]
\(\lambda_p = \pi \sqrt{E/\sigma_p} = \pi \sqrt{206\times10^9 / 200\times10^6} \approx 100.8\)。因 \(\lambda > \lambda_p\)，故欧拉公式适用。
\end{solution}

\begin{example}[title=一端固定一端自由压杆的稳定校核]
圆截面钢压杆，直径 \(d = 40 \,\text{mm}\)，长度 \(l = 800 \,\text{mm}\)，一端固定，一端自由。材料为 Q235，\(\sigma_p = 200 \,\text{MPa}\)，\(E = 206 \,\text{GPa}\)，稳定安全系数 \(n_{\text{st}} = 3\)。若杆顶承受轴向压力 \(F = 30 \,\text{kN}\)，试校核其稳定性。
\end{example}
\begin{solution}
一端固定一端自由：\(\mu = 2\)。
惯性矩：
\[
I = \frac{\pi d^4}{64} = \frac{\pi \times 0.04^4}{64} = 1.257 \times 10^{-7} \,\text{m}^4
\]
惯性半径：
\[
i = \frac{d}{4} = 0.01 \,\text{m},\quad \lambda = \frac{\mu l}{i} = \frac{2 \times 0.8}{0.01} = 160
\]
比例极限柔度 \(\lambda_p \approx 100.8\)，\(\lambda > \lambda_p\)，属大柔度杆，用欧拉公式：
\[
F_{\text{cr}} = \frac{\pi^2 EI}{(\mu l)^2} = \frac{\pi^2 \times 206\times10^9 \times 1.257\times10^{-7}}{(1.6)^2}
= \frac{\pi^2 \times 206 \times 1.257 \times 10^2}{2.56} \approx 1.001 \times 10^5 \,\text{N} = 100.1 \,\text{kN}
\]
工作安全系数：
\[
n = \frac{F_{\text{cr}}}{F} = \frac{100.1}{30} \approx 3.34 > n_{\text{st}} = 3
\]
满足稳定性要求。
\end{solution}

\begin{remark}
压杆稳定性设计时，应尽量使杆在两个形心主惯性平面内的柔度相等（等稳定性），以充分利用材料。例如采用工字形或空心截面时，应使 \(I_{\max}\) 和 \(I_{\min}\) 接近。
\end{remark}